\documentclass[11pt]{article}
\usepackage[margin=1in]{geometry}
\usepackage{amsmath,amssymb}
\usepackage[hidelinks]{hyperref}

\title{Literature-Implied Status Note for arXiv:1406.3842}
\author{run fa\_banach\_001, agent\_lane\_05}
\date{2026-06-14}

\begin{document}
\maketitle

\paragraph{Status.}
\texttt{literature\_implied\_answer} (preprint claim; needs human verification).

\paragraph{Source/open-problem metadata.}
Michael Cwikel, Mario Milman, and Richard Rochberg,
``A brief survey of Nigel Kalton's work on interpolation and related topics,''
arXiv:1406.3842, posted 2014-06-15.
On page 2, Section 2, the paper asks whether one-sided compactness propagates
under interpolation: if an operator is bounded on both endpoints \(A_i\),
\(i=0,1\), and compact on \(A_0\), must it be compact on the interpolation
space \(A_*\)? The paper says that for the complex method the question
``remains unanswered.''

\paragraph{Supporting-answer metadata.}
Evgeniy Pustylnik, ``A new approach to interpolation of compact linear
operators,'' arXiv:2605.00530, posted 2026-05-01. On page 2, Pustylnik
formulates one-sided compactness for embeddings \(A_i\hookrightarrow B_i\),
assumes the interpolation estimate
\[
 \|T\|_{A\to B}\le
 W\bigl(\|T\|_{A_0\to B_0},\|T\|_{A_1\to B_1}\bigr),
\]
where \(W(\alpha,\beta)\to0\) as \(\beta\to0\) uniformly for bounded
\(\alpha\), and explicitly notes that for the complex interpolation functor
\(\mathcal F_{[\theta]}\), \(0<\theta<1\), one may take
\(W(\alpha,\beta)=\alpha^{1-\theta}\beta^\theta\). The Main Theorem on page 2
then asserts compactness of the interpolated embedding. On pages 7--8 the
paper passes from embeddings to arbitrary operators by factoring a continuous
operator \(T:A_0\to B_0\), compact \(T:A_1\to B_1\), through its image spaces.

\paragraph{Identification.}
The 2014 question is the classical one-sided compactness problem for the
complex method. Pustylnik's theorem is stated with compactness at the second
endpoint \(A_1\); the 2014 paragraph phrases compactness at \(A_0\). For the
complex method these are equivalent formulations after swapping the couple
\((A_0,A_1)\) and replacing \(\theta\) by \(1-\theta\). Thus, if the proof of
arXiv:2605.00530 is accepted, it gives a positive answer to the target
question for the complex method.

\paragraph{Provenance limitation.}
The supporting paper does not explicitly cite arXiv:1406.3842 in the checked
source text. This is therefore an agent-identified implication rather than a
paper that explicitly announces it is answering the 2014 survey question. The
supporting source is also a very recent arXiv preprint that claims a broad
resolution of a long-standing problem, so the packet should be human-reviewed
before the run treats the problem as settled.

\paragraph{Search evidence.}
The run indexes were searched for \texttt{1406.3842}, Nigel Kalton,
propagation of compactness, complex method, one-sided compactness, and
Pustylnik/arXiv:2605.00530. No prior packet for this exact target was found.
External web lookup found the arXiv records for arXiv:1406.3842,
arXiv:2109.05968, arXiv:2605.00530, and the related 2014 survey
arXiv:1410.4527. The decisive supporting source used here is arXiv:2605.00530.

\begin{thebibliography}{9}
\bibitem{CMR2014}
M. Cwikel, M. Milman, and R. Rochberg,
``A brief survey of Nigel Kalton's work on interpolation and related topics,''
arXiv:1406.3842, 2014.

\bibitem{Pustylnik2026}
E. Pustylnik,
``A new approach to interpolation of compact linear operators,''
arXiv:2605.00530, 2026.
\end{thebibliography}

\end{document}
